\documentclass[12pt]{article}
\usepackage[utf8]{inputenc}
\usepackage[T1]{fontenc}
\usepackage[margin=1in]{geometry}
\usepackage{setspace}
\usepackage{xcolor}
\usepackage[colorlinks=true, urlcolor=blue!60!black, linkcolor=blue!60!black]{hyperref}
\setlength{\parindent}{0em}
\setlength{\parskip}{0.6em}
\renewcommand{\baselinestretch}{1.0}
\begin{document}
\begin{center}
{\LARGE \textit{Smith v. Doe}, 538 U.S. 84 (2003)\par}
\vspace{0.5em}
{\large Case Brief --- Two-Page Edition\par}
\end{center}

\section{Memory Jogger}
Alaska's retroactive sex offender registration and Internet notification law is civil, not punitive, so it does not violate the Ex Post Facto Clause (\href{https://www.law.cornell.edu/supct/html/01-729.ZO.html}{Cornell LII}).

\section{Detailed Case Facts}
The 1994 Alaska Sex Offender Registration Act required sex offenders, including those convicted earlier, to register and periodically verify personal information, for 15 years or for life, and posted most of it on a public website. Two men convicted of aggravated sexual abuse of a minor before the Act, and released in 1990, challenged its retroactive application.

\section{Procedural History}
The district court ruled for Alaska. The Ninth Circuit found civil intent but punitive effects and reversed. The Supreme Court reversed in 2003.

\section{Judicial Votes}
6--3. Kennedy wrote for Rehnquist, O'Connor, Scalia, and Thomas; Thomas concurred and Souter concurred in the judgment. Stevens dissented, as did Ginsburg, joined by Breyer.

\section{Holding}
The legislature intended a civil regulatory scheme, and respondents failed to show by the ``clearest proof'' that its effects were punitive, so retroactive application is constitutional.

\section{Analysis of Opinions}
\textbf{Majority:} The Act's stated purpose of ``protecting the public from sex offenders'' and its administrative implementation show civil intent. Under the \textit{Mendoza-Martinez} factors, the registry is unlike colonial shaming, imposes no physical restraint, rationally serves public safety, and may treat offenders categorically. \textbf{Thomas:} look only at the statute's face. \textbf{Souter:} the indicators were in ``rough equipoise.'' \textbf{Stevens:} a burden imposed on all and only convicted offenders is punishment. \textbf{Ginsburg:} the Act is excessive because it offers no path off the registry for the rehabilitated.

\section{Examples: Future Applications}
\textbf{Same-Side:} civil commitment of sexually violent predators, as in \textit{Hendricks}. \textbf{Opposite-Side:} applying harsher, later Sentencing Guidelines violates the Clause (\textit{Peugh}). \textbf{Fence-Sitter:} lifetime GPS monitoring supervised by corrections officers.

\section{Critique}
Critics say the ``clearest proof'' standard invites legislatures to relabel punishment as civil, that the ``frightening and high'' recidivism claim rested on a disavowed magazine article, and that the archive analogy ignores modern digital stigma.

\section{Key Quotations}
\begin{quote}
``If the intention of the legislature was to impose punishment, that ends the inquiry. If, however, the intention was to enact a regulatory scheme that is civil and nonpunitive, we must further examine whether the statutory scheme is `so punitive either in purpose or effect as to negate [the State's] intention' to deem it `civil.'\,''
\end{quote}
\noindent --- Justice Kennedy, for the Court, 538 U.S. at 92.

\end{document}
